\documentclass{article}
\usepackage{ctex}
\usepackage[a4paper,
    left=1.75cm,
    right=1.75cm,
    top=2.25cm,
    bottom=2.25cm
]{geometry}
\usepackage{amsmath, amssymb, amsthm}
\usepackage{xcolor}
\usepackage{graphicx}
\usepackage{enumerate}
\usepackage[normalem]{ulem}
\usepackage{subcaption}
\usepackage{fancyhdr}

\newcommand{\highlight}[1]{\textcolor{blue}{#1}}

\pagestyle{fancy}
\fancyhf{} % 清空默认
\fancyfoot[R]{\thepage}

% Problem 环境
\newtheorem{definition}{Definition}
\newtheorem{problem}{Problem}           % 如果要中文可以改成 {问题}
\newtheorem{proposition}{Proposition}
\newtheorem{example}{Example}
\newtheorem{assumption}{Assumption}

\def\highlight#1{\textcolor{blue}{#1}}
\def\hhighlight#1{\textcolor[rgb]{0.765, 0.682, 0.796}{#1}}
\def\corrected#1{\textcolor{cyan}{#1}}

\title{On Cognitive Non-Blocking in Partially Observed Discrete-Event Systems:
Verification and Synthesis using Supervisory Control}
\author{-}
\date{April 2026}

\begin{document}

\maketitle

\section{Introduction}

...

Consider a robotic system deployed in a partially observable environment. The robot has knowledge of the environment; however, it cannot accurately determine its exact current state. 
The objective of the robot is to perform data transmission only when it is located in a safe region. 
Otherwise, critical information may be exposed to potential intruders, leading to security risks.
%
Apparently, in this scenario, safe locations may be observationally indistinguishable from unsafe ones under partial observation. 
To this end, the robot must ensure that data transmission is executed only at locations that can be guaranteed to be safe.

...

\subsection{Contributions and challenges}

The main contributions of this paper are organized as follows.
%
\begin{enumerate}[(i)]
    \item Compared to the existing work of ...

    \item Compared to ....
\end{enumerate}

\subsection{Related works}

\subsection{Organization}

The remainder of this paper is organized as follows.

\section{Preliminaries}\label{sec:prelimi}

This section introduces the preliminaries for the verification and synthesis of cognitive non-blocking in discrete-event systems (DESs), including system models, supervisory control under partial observation, and information states.


\subsection{System model}

\subsection{Supervisory control under partial observation}

We define $\mathcal{L}_m (S/G) := L(S/G) \cap \mathcal{L}(G)$, and
we say that supervisor $S$ is non-blocking if $\overline{\mathcal{L}_m(S/G)} = \mathcal{L}(S/G)$.
One can easily check that the system is non-blocking if and only if 
$\forall s \in \mathcal{L}(S/G), \exists t \in \Sigma^* \ s.t. \ st \in \mathcal{L}_m(S/G)$

...

Given supervisor $S$ and $s \in P(\mathcal{L}(S/G))$ be an observable
string. We define
$$
\mathcal{E}_S(s) := \{x \in X : \exists t \in \mathcal{L}(S/G) \ s.t. \ \delta(t) = x \wedge P(t) = s\}
$$
as the \emph{state estimate} upon the occurrence of $s$, namely, the set of states in which the system could be after observing $s$.
We further define
$$
\hat{\mathcal{E}}_S(s) := \{x \in X : \exists t \in (\Sigma^*\Sigma_o) \cap \mathcal{L}(S/G) s.t. \delta(t) = x \wedge P(t) = s\}
$$
as the \emph{cognitive state estimate} associated with $s$, i.e., the set of states that can be reached immediately after the immediate observation after $s$.

% TODO
% 如果当marking state在不可观区域的话(within an unobservable reach)
% 虽然此时具有到达marking state的能力，但是此时仍然无法确定事情是否被完成
% 通常来说，在典型物理场景中， 如……，能够观测到目标事件被完成，这个过程才算有意义的
In many physical scenarios, such as a robot arriving at a goal region or a machine finishing an operation, completion is expected to be accompanied by some observable event.
However, if a marked state lies within the unobservable reach of a nonmarked state, then the system may enter the marked set silently, and an outside observer cannot determine whether the task has actually been completed.
It is meaningful to interpret the reachability of a marked state as task completion only when such completion can be recognized from the observation.
To this end, we give the following assumption follows.


% TODO
% 跟殷老师讨论，这个不要了，我觉得非常不好
\begin{assumption}
    Given FSA $G=(X,\Sigma,\delta,x_0,X_m)$, we assume that no marked state lies in the unobservable reach of any unmarked state.
    Equivalently, for all states $x \in X$, if its successor $x' \in X_m$ is a marking states, then we have $(x, \sigma_{o}, x') \in \delta$ with $\sigma_o \in \Sigma_o$.
\end{assumption}

\subsection{Information state and update rules}

%
% -----------------------------
\section{Problem formulation}\label{sec:prob}

In this section, we first introduce the notion of cognitive non-blocking based on the preliminaries. 
We then formulate the verification and synthesis problems for DESs under this notion.

\subsection{Cognitive non-blocking}

...

\begin{definition}[Cognitive Non-blocking]
    Given FSA $G = (X, \Sigma, \Delta, q_0, X_m)$ and a set of observable events $\Sigma_o \subseteq \Sigma$, we say FSA $G$ is \textit{cognitive non-blocking} with respect to observable event set $\Sigma_o$ if
    $$
    \forall s \in \mathcal{L}(S/G), \exists t \in \Sigma^* \text{ such that } \hat{\mathcal{E}}_S(st) \subseteq X_m.
    $$
\end{definition}

Intuitively, an FSA satisfies cognitive non-blocking if it can eventually determine, based on its observations, that it is in a marked state \highlight{with certainty}. 
That is, there exists a continuation such that the resulting information state consists entirely of marked states.
It is straightforward to verify that cognitive non-blocking is a stronger property than classical non-blocking.

...

% TODO
\begin{example}
    ...
    
    Although we assume the ...
    It is still difficult to determine whether
    
\end{example}


\subsection{Problem formulation}

\begin{problem}
\textbf{(Cognitive Non-Blocking Verification Problem)}
Given an FSA 
$$
G = (X, \Sigma, \Delta, q_0, Acc)
$$ 
a set of observable events $\Sigma_o \subseteq \Sigma$, and a set of observable events $\Sigma_o \subseteq \Sigma$
verify whether the system $G$ is cognitive non-blocking.
\end{problem}

\begin{problem}
\textbf{(Maximally Permissive Cognitive Non-Blocking Problem)} 
Given an FSA $G$ that is not cognitive non-blocking, a set of observable events $\Sigma_o \subseteq \Sigma$,
and a set of controllable events $\Sigma_c \subseteq \Sigma$, synthesize a partial-observation supervisor $S : P(\mathcal{L}(G)) \rightarrow \Gamma$, such that
\begin{enumerate}[(i)]
    \item the closed-loop system $S / G$ is cognitive non-blocking with respect to $\Sigma_o$;
    \item for any other supervisor $S'$ satisfying (i), we have $\mathcal{L}(S/G) \not \subset \mathcal{L}(S'/G)$.
\end{enumerate}
\end{problem}


\section{Inter-connected system under all-inclusive controller}

In this section, we construct a bipartite transition system (BTS) based on the evolution of information states. 
Then, we introduce the inter-connected system (ICS) that captures the microscopic evolution within the BTS such that the cognitive non-blocking property can be verified.

\subsection{Construction of all-inclusive controller}

We say an all-inclusive is \textit{complete}

\subsection{Inter-Connected System under all-inclusive controller}

\begin{definition}[Inter-Connected System]
    Given an all-inclusive controller $\mathcal{A}$, its corresponding inter-connected system is a tuple
    $\mathcal{C} = (\mathcal{Q}, \Sigma_\mathcal{C}, \delta_\mathcal{C}, q_0, \mathcal{Q}_m)$
    where
\begin{enumerate}
    \item $\mathcal{Q} \subseteq (Y \times X) \cup (Z \times X)$ is the set of states defined by
    \begin{itemize}
        \item $(y, x) \in \mathcal{Q}$ if $y \in Y$ and $x \in y$, and
        \item $(z, x) \in \mathcal{Q}$ if $z \in Z$ and $x \in I(z)$;
    \end{itemize}
    \item $\Sigma_\mathcal{C} = \Sigma \cup \Gamma$ is the set of events;
    \item $\delta_\mathcal{C} \subseteq \mathcal{Q} \times \Sigma_\mathcal{C} \times \mathcal{Q}$ is the transition relation such that for any $\gamma \in \Gamma$,
%     \begin{itemize}
%         \item δICST ((y, x1), γ) = (z, x2) if x1 = x2 and hT Y Z(y, γ) = z
%         \item δICST ((z, x1), σ) = (z, x2) if f(x1, σ) = x2 and σ ∈ Γ(z) ∩ Euo
%         \item δICST ((z, x1), σ) = (y, x2) if f(x1, σ) = x2, σ ∈
% % Γ(z) ∩ Eo and hT ZY (z, σ) = y
%     \end{itemize}
%     \item qICST0 = ({x0}, x0) is the initial state;
%     \item m = {(z, x) ∈ QICST : x ∈ Xm} is the set of
% marked states.

\end{enumerate}
\end{definition}

In the following, we call $...$

\subsection{Region of attraction for cognitive non-blocking}

% 我们认为如果一个Z-state是locally legeal的 其中任意状态均存在到达其中marked Z-microstate的路径，或经由可观事件到达marked Y-microstate的状态转移 接下来，求解该Z状态的SCC 针对这个SCC，通过反向BFS，搜索其中前序Y-状态（不是microstate），最终如果确定这个吸引域能够到达初始Y状态，则认为这个吸引域是【feasible】的 最终通过求解所有吸引域的并集，获得最大允许监督控制器的范围，并得到监督控制器

\begin{definition}
    We say a 
\end{definition}

\subsection{Verification of cognitive non-blocking}

%
%
% -----------------------------
\section{Solution to maximal permissive controller}\label{sec:mpc}

In this section, we construct a maximally permissive controller based on the ICS obtained in Section~\ref{sec:syn_cnb}. 
By leveraging the idea of unfolding livelock Y-states, we transform a memoryless supervisor into a memory-based one, 
such that states trapped in livelock can be guided toward marking states.

\subsection{Supervisor under consistent AIC}

% 这里要提出，尽管supervisor是consistent的，但是仍然不是maximally permissive的

\subsection{Synthesis of cognitive memorable supervisor}

\subsection{Properties of supervisor}

\subsection{Correctness and complexity of the Synthesis Algorithm}


%
% -----------------------------
\section{Conclusion and future works}

\end{document}